% 参考文献与底本原注（p.36）
\section*{参考文献}
\addcontentsline{toc}{section}{参考文献}

\begin{enumerate}[label={[\arabic*]},leftmargin=3.2em,itemsep=0.15em]
\item J. Bernstein. \emph{Harvard lecture notes on representations of $p$-adic reductive groups.}
      （$p$ 进约化群表示论的哈佛讲义。）
\item J. Bernstein, P. Deligne. \emph{Le centre de Bernstein.}
      In: \emph{Représentations des groupes réductifs sur un corps local},
      Travaux en cours (P. Deligne ed.), Hermann, Paris, 1--32 (1984).
      （Bernstein 中心。）
\item J. Bernstein, A. V. Zelevinsky. \emph{Induced representations of reductive $p$-adic groups. I.}
      Ann. Sci. École Norm. Sup. \textbf{10}, 441--472 (1977).
      （约化 $p$ 进群的诱导表示 I。）
\item Bushnell, Kutzko. \emph{The Admissible Dual of $GL(N)$ via Compact Open Subgroups.}
      Annals of Mathematics Studies, 1993.
      （经由紧开子群给出的 $GL(N)$ 的容许对偶。）
\item P. Cartier. \emph{Representations of $p$-adic groups: a survey.}
      Proc. of Symp. in Pure Math. \textbf{33}, part 1, 111--155 (1979).
      （$p$ 进群表示论综述。）
\item W. Casselman. \emph{The unramified principal series of $p$-adic groups I.}
      Compositio Math. \textbf{40}, no. 3, 387--406 (1980).
      （$p$ 进群的非分歧主序列 I。）
\item G. Henniart. \emph{Sur l'unicité des types pour $GL_2$.}
      Duke Mathematical Journal \textbf{2}, no. 115 (2002). Appendix to
      \emph{Multiplicités modulaires et représentations de $GL_2(\mathbb{Z}_p)$
      et de $\mathrm{Gal}(\mathbb{Q}_p/\mathbb{Q}_p)$ en $l=p$}.
      （论 $GL_2$ 类型的唯一性；附录于「模重数与 $l=p$ 时 $GL_2(\mathbb{Z}_p)$
      及 $\mathrm{Gal}(\overline{\mathbb{Q}}_p/\mathbb{Q}_p)$ 的表示」。）
\item M. Harris, R. Taylor. \emph{The geometry and cohomology of some simple Shimura varieties.}
      Annals of Mathematics Studies, vol. 151 (2001).
      （某些简单 Shimura 簇的几何与上同调。）
\item V. Păsku̧nas. \emph{Unicity of types for supercuspidal representations of $GL_N$.}
      Proc. London Math. Soc. (3) \textbf{91}, no. 3, 623--654 (2005).
      （$GL_N$ 超尖点表示类型的唯一性。）
\item J. Tate. \emph{Number theoretic background.}
      In: \emph{Automorphic forms, representations, and $L$-functions}, vol. 2 (1977).
      （数论背景。）
\item T. Wedhorn. \emph{The Local Langlands Correspondence for $GL(n)$ over $p$-adic Fields.}
      Lecture notes (2000).
      （$p$ 进域上 $GL(n)$ 的局部 Langlands 对应。）
\end{enumerate}

\vspace{0.8em}
\noindent\textbf{作者通信地址：}\ \texttt{zijian.yao.math@gmail.com}

\bigskip
\hrule
\bigskip

\section*{原注}
\addcontentsline{toc}{section}{原注}

底本以页脚脚注形式给出 11 条注，正文以上标数字标出。译文照录其位置与内容如下；
凡底本在版式上被截断之处，如实标出，不作补写。

\begin{enumerate}[label={\textsuperscript{\arabic*}},leftmargin=2.4em,itemsep=0.3em]
\item （p.11，$H(G)$ 的定义处）不要与下一小节所称的 Hecke 作用混淆。
\item （p.13）例如，这正是证明 $H(G)=\bigcup_K H(G/\!/K)$ 的方法；或者等价地说，
      两个左、右[不变子空间]……\emph{（底本此注在页脚处被排版截断，末句未完。）}
\item （p.14）在 $K^{n}$ 上引入非退化对称双线性形式之后。
\item （p.14）注意我们其实不能直接应用定理 2.5，因为它的证明用到了 Schur 引理。
\item （p.15）由于时间限制，我作报告时省略了这部分讨论。
\item （p.17）这在 Cartier 的综述 [5] 中称为「对偶 Frobenius 互反」。
\item （p.18）更精确地说，我们应取 $K$-抛物子群 $P\subset G$，然后取其 $K$-点。
\item （p.20）我认为条件 (2)（等价地 (3)）可能正是超尖点性最初的定义方式。
\item （p.23）文献中有时称之为 Bernstein--Zelevinsky 段；本文中我们只用
      「段」或「Zelevinsky 段」。
\item （p.26）更准确地说是 $(\Art^{-1}\chi_1\oplus\Art^{-1}\chi_2,\,0)$。
\item （p.28）二稿中将补充大量内容，尤其是 Whittaker 模型及其在定义
      $\epsilon$[-因子]中的重要性。\emph{（底本此注在页脚处被排版截断。）}
\end{enumerate}
